\section*{Chapter 1 --- Sequences}

\begin{solution}{exo:tle:seq:1}
Let $P(n)$ be the statement $\sum_{k=1}^{n} k^2 = \frac{n(n+1)(2n+1)}{6}$.
\emph{Base case:} for $n = 0$ both sides are $0$ (empty sum).
\emph{Inductive step:} assume $P(n)$. Then
\[
\sum_{k=1}^{n+1} k^2 = \frac{n(n+1)(2n+1)}{6} + (n+1)^2
= \frac{(n+1)\bigl(n(2n+1) + 6(n+1)\bigr)}{6}
= \frac{(n+1)(2n^2 + 7n + 6)}{6}.
\]
Since $2n^2 + 7n + 6 = (n+2)(2n+3)$, this is
$\frac{(n+1)(n+2)(2(n+1)+1)}{6}$, which is $P(n+1)$. By induction, $P(n)$
holds for all $n$.
\end{solution}

\begin{solution}{exo:tle:seq:2}
$a_{n+1} - a_n = \frac{n+2}{n+1} - \frac{n+1}{n}
= \frac{n(n+2) - (n+1)^2}{n(n+1)} = \frac{-1}{n(n+1)} < 0$: $(a_n)$ is
strictly decreasing.

$(b_n)$ has positive terms and
$\frac{b_{n+1}}{b_n} = \frac{2^{n+1}}{n+1}\cdot\frac{n}{2^n} =
\frac{2n}{n+1} \geq 1 \iff 2n \geq n+1 \iff n \geq 1$: $(b_n)$ is increasing
(strictly for $n \geq 2$).

$c_{n+1} - c_n = (n+1)^2 - 10(n+1) - n^2 + 10n = 2n - 9$, which is negative
for $n \leq 4$ and positive for $n \geq 5$: $(c_n)$ decreases up to
$c_5 = -25$, its minimum, then increases. It is not monotonic.
\end{solution}

\begin{solution}{exo:tle:seq:3}
Factor dominant terms:
\[
u_n = \frac{n^2(3 - 1/n + 1/n^2)}{n^2(2 + 5/n^2)} \xrightarrow[n\to+\infty]{} \frac{3}{2}.
\]
Multiply by the conjugate:
\[
v_n = \frac{(n+1) - n}{\sqrt{n+1} + \sqrt{n}} = \frac{1}{\sqrt{n+1}+\sqrt{n}}
\xrightarrow[n\to+\infty]{} 0.
\]
Divide numerator and denominator by $3^n$:
\[
w_n = \frac{(2/3)^n - 1}{1 + (1/3)^n} \xrightarrow[n\to+\infty]{} \frac{0-1}{1+0} = -1,
\]
using $\lim q^n = 0$ for $\abs{q} < 1$.
\end{solution}

\begin{solution}{exo:tle:seq:4}
$u_n = 5 + 3n \to +\infty$ and $v_n = 8 \cdot (1/2)^n = 2^{3-n} \to 0$.
The sums are
\[
\sum_{k=0}^{n} u_k = (n+1)\,\frac{5 + (5+3n)}{2} = \frac{(n+1)(10+3n)}{2}
\xrightarrow[n\to+\infty]{} +\infty,
\]
\[
\sum_{k=0}^{n} v_k = 8\,\frac{1 - (1/2)^{n+1}}{1 - 1/2}
= 16\left(1 - \left(\tfrac{1}{2}\right)^{n+1}\right)
\xrightarrow[n\to+\infty]{} 16 .
\]
\end{solution}

\begin{solution}{exo:tle:seq:5}
Since $-1 \leq \cos n \leq 1$,
\[
\frac{n-1}{n+1} \leq \frac{n + \cos n}{n+1} \leq 1,
\]
and $\frac{n-1}{n+1} \to 1$, so the limit is $1$ by the squeeze theorem.

For the second limit, write
\[
0 \leq \frac{n!}{n^n}
= \frac{1}{n}\cdot\frac{2}{n}\cdots\frac{n}{n}
\leq \frac{1}{n},
\]
because every factor $\frac{k}{n}$ with $2 \leq k \leq n$ is at most $1$.
Since $\frac1n \to 0$, the squeeze theorem gives $\frac{n!}{n^n} \to 0$. (The
suggested geometric bound also works: each factor with $k \leq n/2$ is at most
$\frac12$, giving the stronger bound $(1/2)^{\floor{n/2}}$.)
\end{solution}

\begin{solution}{exo:tle:seq:6}
\emph{1.} $u_0 = 0 \in \intcc{0}{2}$. If $0 \leq u_n \leq 2$, then
$2 \leq u_n + 2 \leq 4$, so $\sqrt{2} \leq u_{n+1} \leq 2$; in particular
$0 \leq u_{n+1} \leq 2$. By induction the property holds for all $n$.

\emph{2.} $u_{n+1} - u_n = \sqrt{u_n + 2} - u_n$. For $x \in \intcc{0}{2}$,
$\sqrt{x+2} \geq x \iff x + 2 \geq x^2 \iff (2-x)(x+1) \geq 0$, which is true.
Hence $(u_n)$ is increasing.

\emph{3.} Increasing and bounded above by $2$, $(u_n)$ converges to some
$\ell \in \intcc{0}{2}$. Passing to the limit in $u_{n+1} = \sqrt{u_n + 2}$
(the map $x \mapsto \sqrt{x+2}$ is continuous) gives $\ell = \sqrt{\ell + 2}$,
so $\ell^2 - \ell - 2 = 0$, i.e.\ $\ell \in \{-1, 2\}$. Since $\ell \geq 0$,
$\lim u_n = 2$.
\end{solution}

\begin{solution}{exo:tle:seq:7}
\emph{1.} Between two doses, $40\%$ of the drug is eliminated, so the
quantity $u_n$ becomes $0.6\,u_n$; the next dose adds $1$ unit:
$u_{n+1} = 0.6\,u_n + 1$.

\emph{2.} $v_{n+1} = u_{n+1} - 2.5 = 0.6\,u_n + 1 - 2.5 = 0.6(u_n - 2.5)
= 0.6\,v_n$: $(v_n)$ is geometric with ratio $0.6$ and first term
$v_0 = 1 - 2.5 = -1.5$. Hence $v_n = -1.5 \times 0.6^n$ and
\[
u_n = 2.5 - 1.5 \times 0.6^n .
\]

\emph{3.} Since $0.6^n \to 0$, $u_n \to 2.5$: the quantity of drug stabilizes
at $2.5$ units.
\end{solution}

\begin{solution}{exo:tle:seq:8}
\emph{1.} $u_0 = 3 > 1$. If $u_n > 1$, then $u_n + 2 > 0$ and
\[
u_{n+1} - 1 = \frac{4u_n - 1 - u_n - 2}{u_n + 2} = \frac{3(u_n - 1)}{u_n + 2} > 0 .
\]
By induction, $u_n > 1$ for all $n$ (and in particular $u_n + 2 \neq 0$, so
the sequence is well defined).

\emph{2.} Using the identity above,
\[
v_{n+1} = \frac{1}{u_{n+1} - 1} = \frac{u_n + 2}{3(u_n - 1)}
= \frac{(u_n - 1) + 3}{3(u_n - 1)} = \frac{1}{3} + v_n .
\]
So $(v_n)$ is arithmetic with common difference $\frac13$ and
$v_0 = \frac{1}{u_0 - 1} = \frac12$.

\emph{3.} $v_n = \frac12 + \frac{n}{3}$, hence
$u_n = 1 + \frac{1}{v_n} = 1 + \frac{6}{3 + 2n}$. Since $v_n \to +\infty$,
$u_n \to 1$.
\end{solution}

\begin{solution}{exo:tle:seq:9}
\emph{1.} $H_{2n} - H_n = \sum_{k=n+1}^{2n} \frac{1}{k}$ is a sum of $n$
terms, each at least $\frac{1}{2n}$; hence
$H_{2n} - H_n \geq n \cdot \frac{1}{2n} = \frac12$.

\emph{2.} By induction on $k$: $H_{2^0} = H_1 = 1 \geq 1$. If
$H_{2^k} \geq 1 + \frac{k}{2}$, then applying point 1 with $n = 2^k$,
\[
H_{2^{k+1}} \geq H_{2^k} + \frac12 \geq 1 + \frac{k+1}{2}.
\]
The sequence $(H_n)$ is increasing (each step adds $\frac{1}{n+1} > 0$) and
the subsequence $H_{2^k}$ is unbounded, so $(H_n)$ is not bounded above.
Increasing and unbounded, it tends to $+\infty$
(\cref{thm:tle:seq:monotone}).
\end{solution}

\begin{solution}{exo:tle:seq:10}
\emph{1.} The sequence $d_n = b_n - a_n$ satisfies
$d_{n+1} - d_n = (b_{n+1} - b_n) - (a_{n+1} - a_n) \leq 0$, so $(d_n)$ is
decreasing; since $d_n \to 0$, we get $d_n \geq 0$ for all $n$ (a decreasing
sequence with a negative term would stay below it forever, preventing the
limit $0$). Hence $a_n \leq b_n$.

\emph{2.} From $a_n \leq b_n \leq b_0$, the increasing sequence $(a_n)$ is
bounded above, so it converges to some $\ell$. Similarly $(b_n)$, decreasing
and bounded below by $a_0$, converges to some $\ell'$. Then
$\ell' - \ell = \lim (b_n - a_n) = 0$, so $\ell = \ell'$.

\emph{3.} $(a_n)$ is (strictly) increasing since
$a_{n+1} - a_n = \frac{1}{(n+1)!} > 0$. For $(b_n)$,
\[
b_{n+1} - b_n = \frac{1}{(n+1)!} + \frac{1}{(n+1)(n+1)!} - \frac{1}{n\,n!}
= \frac{n(n+1) + n - (n+1)^2}{n(n+1)(n+1)!}
= \frac{-1}{n(n+1)(n+1)!} < 0 ,
\]
so $(b_n)$ is decreasing. Finally $b_n - a_n = \frac{1}{n\,n!} \to 0$. The
two sequences are adjacent, hence converge to a common limit.
\end{solution}
